% Generated draft: requires mathematical and source review.
\documentclass[11pt]{article}
\usepackage[margin=25mm]{geometry}
\usepackage{fontspec}
\setmainfont{Latin Modern Roman}
\usepackage{amsmath,amssymb,mathtools,unicode-math}
\setmathfont{Latin Modern Math}
\usepackage{enumitem,needspace,xurl,hyperref}
\setlength{\parindent}{0pt}
\setlength{\parskip}{5pt}
\setlength{\emergencystretch}{3em}
\newcommand{\N}{\mathbb{N}}
\newcommand{\Z}{\mathbb{Z}}
\newcommand{\Q}{\mathbb{Q}}
\newcommand{\R}{\mathbb{R}}
\newcommand{\C}{\mathbb{C}}
\newcommand{\F}{\mathbb{F}}
\newcommand{\A}{\mathbb{A}}
\newcommand{\PP}{\mathbb{P}}
\newcommand{\E}{\mathbb{E}}
\newcommand{\eps}{\varepsilon}
\newcommand{\dd}{\,\mathrm{d}}
\begin{document}
\section*{UnsolvedMath ID: 700452}
\textbf{Source problem code:} catalogue entry
\subsection{Definitions and mathematical statement}
Let \(F\) be an absolutely free group with a chosen free basis \(X=\{x_i:i\in I\}\), let \(R\trianglelefteq F\), and let \(\mathfrak V\) be any variety of groups. No finiteness assumption on \(I\) or \(F/R\) is imposed. For a group \(H\), let \(\mathfrak V(H)\) denote the verbal subgroup generated by all values in \(H\) of words that are identities in \(\mathfrak V\). Put \(Q=F/R\) and \(q_i=x_iR\).

A relatively free group in \(\mathfrak V\) on a specified set \(S\) means the quotient of the absolutely free group on \(S\) by its \(\mathfrak V\)-verbal subgroup. Let \(A\) be relatively free in \(\mathfrak V\) on generators indexed by \(I\). Use the following basis-dependent realization of its \(\mathfrak V\)-verbal wreath product with \(Q\). Let \(B\) be relatively free in \(\mathfrak V\) on the symbols
\[
\{a_{i,t}:i\in I,\ t\in Q\}.
\]
Let \(Q\) act on \(B\) by \(q\cdot a_{i,t}=a_{i,qt}\), and set
\[
W=B\rtimes Q,\qquad (b,q)(c,r)=(b(q\cdot c),qr).
\]
The embedding recalled in the source is represented, with these conventions, by the induced map
\[
\iota:F/\mathfrak V(R)\longrightarrow W,
\qquad
\iota\bigl(x_i\mathfrak V(R)\bigr)=(a_{i,1_Q},q_i).
\]
The source supplies the assertion that this induced map is an isomorphic embedding.

\textbf{Problem.} For arbitrary \(F\), \(R\), and \(\mathfrak V\) as above, find a criterion characterizing precisely the elements of \(W\) in the image of \(\iota\); that is, give necessary and sufficient conditions on \((b,q)\in B\rtimes Q\) for
\[
(b,q)\in\iota\bigl(F/\mathfrak V(R)\bigr).
\]
The source does not explicitly specify the embedding or wreath-product conventions; the realization above fixes these choices for this draft and requires comparison with the cited construction during mathematical review. The source states that a criterion is already known when \(\mathfrak V\) is the variety of all abelian groups. It also records an explicitly unverified AI claim of a general criterion; that claim is not adopted here as an established solution.
\subsection{Short English statement}
Given a free group \(F\), a normal subgroup \(R\), and a variety of groups \(\mathfrak V\), characterize the image of the source's embedding of \(F/\mathfrak V(R)\) into the corresponding verbal wreath product. The source reports a known criterion for the abelian variety and an unverified AI claim for the general problem.
\subsection{Sources}
\begin{itemize}
\item Local UnsolvedMath registry: permanent ID \(700452\), exact source problem code \texttt{catalogue entry}, title \emph{Kourovka Notebook Problem 7.58}; supplied upstream revision \texttt{d371ee917ac9647fbb06cd6fc529deec6f3dda82}.
\item G. G. Yabanzhi, proposer, 1980. \emph{The Kourovka Notebook (Unsolved Problems in Group Theory)}, problem 7.58, originating in the 7th issue (1980); supplied record references the 21st issue edition (2026), \texttt{arXiv:1401.0300v48}. The supplied collection description names E. I. Khukhro and V. D. Mazurov as editors.
\item H. Neumann, \emph{Varieties of Groups}, Springer, Berlin, 1967. Cited by the supplied statement for the embedding assertion.
\item V. N. Remeslennikov and V. G. Sokolov, \emph{Algebra and Logic}, \textbf{9}, no. 5 (1970), 342--349. Cited by the supplied statement for the criterion in the abelian case; no article title is supplied.
\item Supplied discussion: ``[Unverified AI claim] A criterion is found,'' attributed to Achyuth Jayadevan, \emph{Preprint}, 2026, \url{https://zenodo.org/records/22821158}. The supplied status review says that the 21st issue (October 2026), \texttt{arXiv:1401.0300v48}, marks this claim with \(\#\) as an unverified AI solution claim.
\item Supplied dated curation: status review checked at \texttt{2026-10-16}, with canonical status \texttt{open} and evidence label \texttt{unrefereed\_claim}. This is supplied dated curation, not an independently verified literature assessment. The record's status label is also \texttt{open}.
\end{itemize}
\end{document}
